\section{Mitigación aplicada por Meta}

\textbf{Medidas inmediatas.} Como no había acceso remoto, los ingenieros fueron en persona a los centros de datos, depuraron y reiniciaron los equipos y volvieron a levantar la red troncal \cite{meta_detalles}. Para devolver el servicio no encendieron todo a la vez. Algunos centros de datos habían bajado su consumo eléctrico en decenas de megavatios, y subirlo de golpe podía dañar los sistemas eléctricos y los cachés. Por eso devolvieron el tráfico por etapas, con el procedimiento que practican en sus simulacros \emph{storm}, donde apagan un servicio, un centro de datos o una región entera \cite{meta_detalles}.

\textbf{Cambios permanentes.} Meta anunció que reforzaría sus pruebas, sus simulacros y su resiliencia en general \cite{meta_detalles}. Sus simulacros nunca habían probado la caída de toda la red troncal \cite{meta_detalles}. La empresa decidió mantener la seguridad física estricta de sus centros de datos, aunque eso hiciera más lenta la recuperación, porque una falla así debería ser poco frecuente \cite{meta_detalles}.
